\IfFileExists{stacks-project.cls}{%
\documentclass{stacks-project}
}{%
\documentclass{amsart}
}

% For dealing with references we use the comment environment
\usepackage{verbatim}
\newenvironment{reference}{\comment}{\endcomment}
%\newenvironment{reference}{}{}
\newenvironment{slogan}{\comment}{\endcomment}
\newenvironment{history}{\comment}{\endcomment}

% For commutative diagrams we use Xy-pic
\usepackage[all]{xy}

% We use 2cell for 2-commutative diagrams.
\xyoption{2cell}
\UseAllTwocells

% We use multicol for the list of chapters between chapters
\usepackage{multicol}

% This is generally recommended for better output
\usepackage{lmodern}
\usepackage[T1]{fontenc}

% For cross-file-references
\usepackage{xr-hyper}

% Package for hypertext links:
\usepackage{hyperref}

% For any local file, say "hello.tex" you want to link to please
% use \externaldocument[hello-]{hello}
\externaldocument[introduction-]{introduction}
\externaldocument[conventions-]{conventions}
\externaldocument[sets-]{sets}
\externaldocument[categories-]{categories}
\externaldocument[topology-]{topology}
\externaldocument[sheaves-]{sheaves}
\externaldocument[sites-]{sites}
\externaldocument[stacks-]{stacks}
\externaldocument[fields-]{fields}
\externaldocument[algebra-]{algebra}
\externaldocument[brauer-]{brauer}
\externaldocument[homology-]{homology}
\externaldocument[derived-]{derived}
\externaldocument[simplicial-]{simplicial}
\externaldocument[more-algebra-]{more-algebra}
\externaldocument[smoothing-]{smoothing}
\externaldocument[modules-]{modules}
\externaldocument[sites-modules-]{sites-modules}
\externaldocument[injectives-]{injectives}
\externaldocument[cohomology-]{cohomology}
\externaldocument[sites-cohomology-]{sites-cohomology}
\externaldocument[dga-]{dga}
\externaldocument[dpa-]{dpa}
\externaldocument[sdga-]{sdga}
\externaldocument[hypercovering-]{hypercovering}
\externaldocument[schemes-]{schemes}
\externaldocument[constructions-]{constructions}
\externaldocument[properties-]{properties}
\externaldocument[morphisms-]{morphisms}
\externaldocument[coherent-]{coherent}
\externaldocument[divisors-]{divisors}
\externaldocument[limits-]{limits}
\externaldocument[varieties-]{varieties}
\externaldocument[topologies-]{topologies}
\externaldocument[descent-]{descent}
\externaldocument[perfect-]{perfect}
\externaldocument[more-morphisms-]{more-morphisms}
\externaldocument[flat-]{flat}
\externaldocument[groupoids-]{groupoids}
\externaldocument[more-groupoids-]{more-groupoids}
\externaldocument[etale-]{etale}
\externaldocument[chow-]{chow}
\externaldocument[intersection-]{intersection}
\externaldocument[pic-]{pic}
\externaldocument[weil-]{weil}
\externaldocument[adequate-]{adequate}
\externaldocument[dualizing-]{dualizing}
\externaldocument[duality-]{duality}
\externaldocument[discriminant-]{discriminant}
\externaldocument[derham-]{derham}
\externaldocument[local-cohomology-]{local-cohomology}
\externaldocument[algebraization-]{algebraization}
\externaldocument[curves-]{curves}
\externaldocument[resolve-]{resolve}
\externaldocument[models-]{models}
\externaldocument[functors-]{functors}
\externaldocument[equiv-]{equiv}
\externaldocument[pione-]{pione}
\externaldocument[etale-cohomology-]{etale-cohomology}
\externaldocument[proetale-]{proetale}
\externaldocument[relative-cycles-]{relative-cycles}
\externaldocument[more-etale-]{more-etale}
\externaldocument[trace-]{trace}
\externaldocument[crystalline-]{crystalline}
\externaldocument[spaces-]{spaces}
\externaldocument[spaces-properties-]{spaces-properties}
\externaldocument[spaces-morphisms-]{spaces-morphisms}
\externaldocument[decent-spaces-]{decent-spaces}
\externaldocument[spaces-cohomology-]{spaces-cohomology}
\externaldocument[spaces-limits-]{spaces-limits}
\externaldocument[spaces-divisors-]{spaces-divisors}
\externaldocument[spaces-over-fields-]{spaces-over-fields}
\externaldocument[spaces-topologies-]{spaces-topologies}
\externaldocument[spaces-descent-]{spaces-descent}
\externaldocument[spaces-perfect-]{spaces-perfect}
\externaldocument[spaces-more-morphisms-]{spaces-more-morphisms}
\externaldocument[spaces-flat-]{spaces-flat}
\externaldocument[spaces-groupoids-]{spaces-groupoids}
\externaldocument[spaces-more-groupoids-]{spaces-more-groupoids}
\externaldocument[bootstrap-]{bootstrap}
\externaldocument[spaces-pushouts-]{spaces-pushouts}
\externaldocument[spaces-chow-]{spaces-chow}
\externaldocument[groupoids-quotients-]{groupoids-quotients}
\externaldocument[spaces-more-cohomology-]{spaces-more-cohomology}
\externaldocument[spaces-simplicial-]{spaces-simplicial}
\externaldocument[spaces-duality-]{spaces-duality}
\externaldocument[formal-spaces-]{formal-spaces}
\externaldocument[restricted-]{restricted}
\externaldocument[spaces-resolve-]{spaces-resolve}
\externaldocument[formal-defos-]{formal-defos}
\externaldocument[defos-]{defos}
\externaldocument[cotangent-]{cotangent}
\externaldocument[examples-defos-]{examples-defos}
\externaldocument[algebraic-]{algebraic}
\externaldocument[examples-stacks-]{examples-stacks}
\externaldocument[stacks-sheaves-]{stacks-sheaves}
\externaldocument[criteria-]{criteria}
\externaldocument[artin-]{artin}
\externaldocument[quot-]{quot}
\externaldocument[stacks-properties-]{stacks-properties}
\externaldocument[stacks-morphisms-]{stacks-morphisms}
\externaldocument[stacks-limits-]{stacks-limits}
\externaldocument[stacks-cohomology-]{stacks-cohomology}
\externaldocument[stacks-perfect-]{stacks-perfect}
\externaldocument[stacks-introduction-]{stacks-introduction}
\externaldocument[stacks-more-morphisms-]{stacks-more-morphisms}
\externaldocument[stacks-geometry-]{stacks-geometry}
\externaldocument[moduli-]{moduli}
\externaldocument[moduli-curves-]{moduli-curves}
\externaldocument[examples-]{examples}
\externaldocument[exercises-]{exercises}
\externaldocument[guide-]{guide}
\externaldocument[desirables-]{desirables}
\externaldocument[coding-]{coding}
\externaldocument[obsolete-]{obsolete}
\externaldocument[fdl-]{fdl}
\externaldocument[index-]{index}

% Theorem environments.
%
\theoremstyle{plain}
\newtheorem{theorem}[subsection]{Theorem}
\newtheorem{proposition}[subsection]{Proposition}
\newtheorem{lemma}[subsection]{Lemma}

\theoremstyle{definition}
\newtheorem{definition}[subsection]{Definition}
\newtheorem{example}[subsection]{Example}
\newtheorem{exercise}[subsection]{Exercise}
\newtheorem{situation}[subsection]{Situation}

\theoremstyle{remark}
\newtheorem{remark}[subsection]{Remark}
\newtheorem{remarks}[subsection]{Remarks}

\numberwithin{equation}{subsection}

% Macros
%
\def\lim{\mathop{\mathrm{lim}}\nolimits}
\def\colim{\mathop{\mathrm{colim}}\nolimits}
\def\Spec{\mathop{\mathrm{Spec}}}
\def\Hom{\mathop{\mathrm{Hom}}\nolimits}
\def\Ext{\mathop{\mathrm{Ext}}\nolimits}
\def\SheafHom{\mathop{\mathcal{H}\!\mathit{om}}\nolimits}
\def\SheafExt{\mathop{\mathcal{E}\!\mathit{xt}}\nolimits}
\def\Sch{\mathit{Sch}}
\def\Mor{\mathop{\mathrm{Mor}}\nolimits}
\def\Ob{\mathop{\mathrm{Ob}}\nolimits}
\def\Sh{\mathop{\mathit{Sh}}\nolimits}
\def\NL{\mathop{N\!L}\nolimits}
\def\CH{\mathop{\mathrm{CH}}\nolimits}
\def\proetale{{pro\text{-}\acute{e}tale}}
\def\etale{{\acute{e}tale}}
\def\QCoh{\mathit{QCoh}}
\def\Ker{\mathop{\mathrm{Ker}}}
\def\Im{\mathop{\mathrm{Im}}}
\def\Coker{\mathop{\mathrm{Coker}}}
\def\Coim{\mathop{\mathrm{Coim}}}

% Boxtimes
%
\DeclareMathSymbol{\boxtimes}{\mathbin}{AMSa}{"02}

%
% Macros for moduli stacks/spaces
%
\def\QCohstack{\mathcal{QC}\!\mathit{oh}}
\def\Cohstack{\mathcal{C}\!\mathit{oh}}
\def\Spacesstack{\mathcal{S}\!\mathit{paces}}
\def\Quotfunctor{\mathrm{Quot}}
\def\Hilbfunctor{\mathrm{Hilb}}
\def\Curvesstack{\mathcal{C}\!\mathit{urves}}
\def\Polarizedstack{\mathcal{P}\!\mathit{olarized}}
\def\Complexesstack{\mathcal{C}\!\mathit{omplexes}}
% \Pic is the operator that assigns to X its picard group, usage \Pic(X)
% \Picardstack_{X/B} denotes the Picard stack of X over B
% \Picardfunctor_{X/B} denotes the Picard functor of X over B
\def\Pic{\mathop{\mathrm{Pic}}\nolimits}
\def\Picardstack{\mathcal{P}\!\mathit{ic}}
\def\Picardfunctor{\mathrm{Pic}}
\def\Deformationcategory{\mathcal{D}\!\mathit{ef}}
\usepackage{cleveref}
\crefname{originallemma}{original lemma}{original lemmas}
\Crefname{originallemma}{Original lemma}{Original lemmas}
\makeatletter
\newlabel{lemma-groupoid-on-field-homogeneous}{{04LP}{source}{}{sourcepointer-1}{}}
\newlabel{lemma-groupoid-on-field-homogeneous@cref}{{[originallemma][0][]04LP}{[1][0][]source}{}{}{}}
\makeatother
\crefname{originallemma}{original lemma}{original lemmas}
\Crefname{originallemma}{Original lemma}{Original lemmas}
\makeatletter
\newlabel{morphisms-lemma-dimension-fibre-after-base-change}{{02FY}{source}{}{sourcepointer-2}{}}
\newlabel{morphisms-lemma-dimension-fibre-after-base-change@cref}{{[originallemma][0][]02FY}{[1][0][]source}{}{}{}}
\makeatother
\crefname{originallemma}{original lemma}{original lemmas}
\Crefname{originallemma}{Original lemma}{Original lemmas}
\makeatletter
\newlabel{morphisms-lemma-dimension-fibre-at-a-point}{{02FX}{source}{}{sourcepointer-3}{}}
\newlabel{morphisms-lemma-dimension-fibre-at-a-point@cref}{{[originallemma][0][]02FX}{[1][0][]source}{}{}{}}
\makeatother
\crefname{originalreference}{original reference}{original references}
\Crefname{originalreference}{Original reference}{Original references}
\makeatletter
\newlabel{morphisms-section-points-finite-type}{{01T9}{source}{}{sourcepointer-4}{}}
\newlabel{morphisms-section-points-finite-type@cref}{{[originalreference][0][]01T9}{[1][0][]source}{}{}{}}
\makeatother
\crefname{originallemma}{original lemma}{original lemmas}
\Crefname{originallemma}{Original lemma}{Original lemmas}
\makeatletter
\newlabel{properties-lemma-codimension-local-ring}{{02IZ}{source}{}{sourcepointer-5}{}}
\newlabel{properties-lemma-codimension-local-ring@cref}{{[originallemma][0][]02IZ}{[1][0][]source}{}{}{}}
\makeatother
\title{Dimensions of a groupoid scheme over a field with s, t locally of finite type}
\author{The Stacks Project Authors}
\date{stacks-project:more-groupoids.tex:local-snapshot-20261012; GNU FDL 1.2 or later}
\begin{document}
\maketitle
\setcounter{subsection}{1}
\noindent\textit{Editorial note.} All delivered mathematical text is the author's own, copied from the pinned Stacks Project chapter file more-groupoids.tex; no statement or proof step was rewritten, reordered or removed. The counted span is the author environment and its own proof environment. References inside the span are handed to the assembler as navigation stubs that carry the canonical tag of the cited tag only; whether each of them is load-bearing is recorded on the stub itself.\par\medskip
\noindent\textit{Editorial note.} Manual re-read applied under the pre-screen escalation rule (proof\_bytes 814, at or above the 800-byte threshold): all 22 proof lines re-read in full; each of the four conclusions carries an explicit derivation in the author text; no details-omitted remarks, no sketches, no as-in-the-proof-of redirections; judged complete and accepted.\par\medskip
\noindent\textit{Allowed standard prerequisites (editorial).}\begin{itemize}
\item Transitivity of groupoid arrows over a field: for any r, r' in R a field extension k'/k and an isomorphism of the s-fibres over Spec(k') carrying a point over r to a point over r' (more-groupoids chapter, Lemma lemma-groupoid-on-field-homogeneous, canonical tag 04LP; also in this lane's queue). cited at proof lines 1531-1535. Equalizes the local dimensions dim\_r(R) and dim\_\{r'\}(R) after transport over the extension.
\item Invariance of the dimension of a fibre under extension of the base field (Morphisms chapter, Lemma morphisms-lemma-dimension-fibre-after-base-change, canonical tag 02FY). cited at proof lines 1532-1535. Ensures the local dimensions computed over Spec(k') equal the original ones over Spec(k), so the transport gives dim\_r(R) = dim\_\{r'\}(R) over k itself.
\item Dimension formula at a point: dim\_x(X) = dim(O\_\{X,x\}) + trdeg of the residue field over the residue field of the image point for a morphism locally of finite type (Morphisms chapter, Lemma morphisms-lemma-dimension-fibre-at-a-point, canonical tag 02FX). cited at proof lines 1536-1543. Applied with s and with t it gives the displayed chain dim\_r(R) = dim(O\_\{R,r\}) + trdeg\_\{s(k)\}(kappa(r)) = dim(O\_\{R,r\}) + trdeg\_\{t(k)\}(kappa(r)), the core identity behind items (2), (3) and (4).
\item Hilbert Nullstellensatz: closed points of a scheme locally of finite type over a field k have residue field finite over k (Morphisms chapter, Section morphisms-section-points-finite-type, canonical tag 01T9; kind reference). cited at proof lines 1548-1549. Gives trdeg\_k(kappa(r)) = 0 for closed points r, the input that pins the common value d at closed points.
\item Dimension of a Noetherian scheme as supremum of local dimensions (Properties chapter, Lemma properties-lemma-codimension-local-ring, canonical tag 02IZ). cited at proof lines 1544-1546. Identifies dim(R), and the dimension of any open subset of R, with the supremum of dim(O\_\{R,r\}) over points r, which is maximal at closed points.
\end{itemize}\medskip
\ifdefined\hypertarget\hypertarget{sourcepointer-1}{}\fi
\noindent\textit{Source reference.} Original lemma 04LP: sources/stacks-chow/tags.txt, source lines 6015--6015. Cited at proof lines 1531-1535: for any r, r' it provides a field extension and an isomorphism of the s-fibres over Spec(k') carrying r to r', used to equalize dim\_r(R) and dim\_\{r'\}(R).\par\medskip
\ifdefined\hypertarget\hypertarget{sourcepointer-2}{}\fi
\noindent\textit{Source reference.} Original lemma 02FY: sources/stacks-chow/tags.txt, source lines 3209--3209. Cited at proof lines 1532-1535: the dimension of a fibre is unchanged by extension of the base field, so the transported local dimensions over Spec(k') equal the original ones.\par\medskip
\ifdefined\hypertarget\hypertarget{sourcepointer-3}{}\fi
\noindent\textit{Source reference.} Original lemma 02FX: sources/stacks-chow/tags.txt, source lines 3208--3208. Cited at proof lines 1536-1543: applied with s and with t it yields the chain dim\_r(R) = dim(O\_\{R,r\}) + trdeg\_\{s(k)\}(kappa(r)) = dim(O\_\{R,r\}) + trdeg\_\{t(k)\}(kappa(r)).\par\medskip
\ifdefined\hypertarget\hypertarget{sourcepointer-4}{}\fi
\noindent\textit{Source reference.} Original reference 01T9: sources/stacks-chow/tags.txt, source lines 2366--2366. Cited at proof lines 1548-1549: the Hilbert Nullstellensatz material in this section gives trdeg\_k(kappa(r)) = 0 for closed points r of the finite type scheme R.\par\medskip
\ifdefined\hypertarget\hypertarget{sourcepointer-5}{}\fi
\noindent\textit{Source reference.} Original lemma 02IZ: sources/stacks-chow/tags.txt, source lines 3318--3318. Cited at proof lines 1544-1546: dim(R), and the dimension of any open subset of R, is the supremum of the dimensions of the local rings of R.\par\medskip
\begin{lemma}
\label{lemma-groupoid-on-field-locally-finite-type-dimension}
Let $S$ be a scheme. Let $(U, R, s, t, c)$ be a groupoid scheme
over $S$. Assume $U = \Spec(k)$ with $k$ a field.
Assume $s, t$ are locally of finite type.
Then
\begin{enumerate}
\item $R$ is equidimensional,
\item $\dim(R) = \dim_r(R)$ for all $r \in R$,
\item for any $r \in R$ we have
$\text{trdeg}_{s(k)}(\kappa(r)) = \text{trdeg}_{t(k)}(\kappa(r))$, and
\item for any closed point $r \in R$ we have
$\dim(R) = \dim(\mathcal{O}_{R, r})$.
\end{enumerate}
\end{lemma}

\begin{proof}
Let $r, r' \in R$.
Then $\dim_r(R) = \dim_{r'}(R)$ by
Lemma \ref{lemma-groupoid-on-field-homogeneous}
and
Morphisms, Lemma \ref{morphisms-lemma-dimension-fibre-after-base-change}.
By
Morphisms, Lemma \ref{morphisms-lemma-dimension-fibre-at-a-point}
we have
$$
\dim_r(R) =
\dim(\mathcal{O}_{R, r}) + \text{trdeg}_{s(k)}(\kappa(r)) =
\dim(\mathcal{O}_{R, r}) + \text{trdeg}_{t(k)}(\kappa(r)).
$$
On the other hand, the dimension of $R$ (or any open subset of $R$)
is the supremum of the dimensions of the local rings of $R$, see
Properties, Lemma \ref{properties-lemma-codimension-local-ring}.
Clearly this is maximal for closed points $r$ in which case
$\text{trdeg}_k(\kappa(r)) = 0$ (by the Hilbert Nullstellensatz, see
Morphisms, Section \ref{morphisms-section-points-finite-type}).
Hence the lemma follows.
\end{proof}

\end{document}
